\section{ماژول hybrid\_svd\_layer (Hybrid SVD Layer)}
\label{sec:module_hybrid_svd_layer}

\subsection{معرفی و انگیزه}
این مکانیزم بر مبنای مبانی نظری پیشرفته در ادبیات توجه و فشرده‌سازی \cite{ref_hybridsvdlayer_2024} توسعه یافته است.

تجزیه مقادیر منفرد (\lr{SVD}) در معماری‌های توجه دو رویکرد بنیادین دارد: فشرده‌سازی بر روی فعال‌سازی‌های برخط لایه توجه (\lr{Activation-based SVD Attention}) و فشرده‌سازی ایستای ماتریس‌های وزن لایه‌های خطی (\lr{Weight-based SVD Compression}). رویکرد نخست که در نسخه‌های اولیه تحت عنوان \lr{SVDSingleHead} پیاده‌سازی شده بود، \lr{SVD} را به صورت بلادرنگ بر روی تانسور فعال‌سازی ورودی اعمال می‌کرد. با این حال، محاسبه برخط تجزیه مقادیر منفرد بر روی فعال‌سازی‌ها دارای مرتبه پیچیدگی زمانی سنگین $\mathcal{O}(N^2 d)$ بوده و از نظر پایداری گرادیان در طول فاز پسرو به شدت ناپایدار و از نظر تئوری نامناسب است.

موتور نوین فشرده‌سازی \lr{Hybrid SVD Layer} با بازطراحی کامل ساختار قبلی، فشرده‌سازی ایستای مبتنی بر وزن را با معماری انکودر ترکیب نموده است. این ماژول ماتریس‌های پروجکشن توجه و شبکه‌های پیش‌خوران را با لایه‌های کم‌مرتبه \lr{SVDLinear} جایگزین می‌نماید و با پشتیبانی هم‌زمان از ساختار ماژول خطی فاکتورگیری‌شده، سیستم پایش خطای فروبنیوس و مسیر بازگشت به حالت متراکم اولیه، کاهش پایدار حجم پارامترها و تسریع زمان استنتاج را محقق می‌سازد.

\subsection{مبانی نظری و فرمول‌بندی ریاضی}
در لایه فاکتورگیری‌شده مبتنی بر وزن، ماتریس وزن $\mathbf{W} \in \mathbb{R}^{d_{\text{out}} \times d_{\text{in}}}$ با رتبه انتخابی $r \le \min(d_{\text{in}}, d_{\text{out}})$ از طریق تجزیه مقادیر منفرد اقتصادی (\lr{Thin SVD}) تجزیه می‌شود:
\begin{equation}
\mathbf{W} = \mathbf{U} \mathbf{\Sigma} \mathbf{V}^\top \approx \mathbf{U}_r \mathbf{\Sigma}_r \mathbf{V}_r^\top
\end{equation}
با توزیع متقارن مقادیر منفرد بر روی مولفه‌های چپ و راست، دو ماتریس تجزیه‌شده به دست می‌آیند:
\begin{equation}
\mathbf{W}_1 = \mathbf{\Sigma}_r^{1/2} \mathbf{V}_r^\top \in \mathbb{R}^{r \times d_{\text{in}}}, \quad \mathbf{W}_2 = \mathbf{U}_r \mathbf{\Sigma}_r^{1/2} \in \mathbb{R}^{d_{\text{out}} \times r}
\end{equation}
گذر پیشرو برای تانسور ورودی $\mathbf{x} \in \mathbb{R}^{B \times N \times d_{\text{in}}}$ به دو ضرب متوالی ماتریسی شکسته می‌شود:
\begin{equation}
\mathbf{y} = (\mathbf{x} \mathbf{W}_1^\top) \mathbf{W}_2^\top + \mathbf{b} + \frac{\alpha}{r_{\text{lora}}} (\mathbf{x} \mathbf{A}^\top) \mathbf{B}^\top
\end{equation}
خطای نسبی بازسازی فروبنیوس در هر لایه محاسبه شده و با آستانه مجاز $\tau_{\text{max}}$ مقایسه می‌گردد:
\begin{equation}
\mathcal{E}_{\text{rel}} = \frac{\|\mathbf{W} - \mathbf{W}_2 \mathbf{W}_1\|_F}{\|\mathbf{W}\|_F}
\end{equation}
اگر $\mathcal{E}_{\text{rel}} > \tau_{\text{max}}$، فرآیند فشرده‌سازی لایه لغو شده و ساختار اصلی لایه دست‌نخورده حفظ می‌گردد.

\paragraph{تحلیل طیفی و نامساوی‌های پایداری خطا}
برای بررسی میزان دقت تقریب ماتریس وزن یا نگاشت توجه در ماژول \lr{hybrid\_svd\_layer}، تفاضل ماتریس فاکتورگیری‌شده از نگاشت متراکم اصلی بر اساس نامساوی هافدینگ-فون‌نیومان ارزیابی می‌گردد. فرض کنید $\mathbf{W} \in \mathbb{R}^{d_{\text{in}} \times d_{\text{out}}}$ ماتریس وزن هدف باشد که توسط تقریب کم‌مرتبه یا پراکنده با رتبه موثر $r$ جایگزین شده است. کران بالای خطای بازسازی بر مبنای نرم فروبنیوس به صورت زیر تصریح می‌شود:
\begin{equation}
\label{eq:spectral_bound_hybrid_svd_layer}
\| \mathbf{W} - \mathbf{W}_{r} \|_F = \sqrt{\sum_{i=r+1}^{\min(d_{\text{in}}, d_{\text{out}})} \sigma_i^2(\mathbf{W})} \le \epsilon \| \mathbf{W} \|_F
\end{equation}
که در آن $\sigma_i(\mathbf{W})$ بیانگر $i$-امین مقدار منفرد مرتب‌شده ماتریس به صورت نزولی است. نسبت حفظ انرژی طیفی $\eta_r = \frac{\sum_{i=1}^r \sigma_i^2}{\sum_{i=1}^{d} \sigma_i^2}$ تضمین می‌کند که با انتخاب آستانه $\tau \ge 0.90$، بیش از نود درصد از پراکندگی واریانس داده‌ها حفظ شده و افت سیگنال در انتشار گرادیان رو به عقب به حداقل ممکن کاهش می‌یابد.

\paragraph{تحلیل مجانبی و پیچیدگی اشتراک حافظه سخت‌افزاری}
از منظر پیچیدگی محاسباتی و انتقال داده‌ها در سلسومراتب حافظه (حافظه سراسری \lr{HBM} به حافظه محلی \lr{SRAM})، عملیات ضرب ماتریسی در این مکانیزم نیازمند تخصیص بافر موقت با ابعاد فاکتورگیری‌شده است. برای پردازش دنباله‌ای به طول $N$ با بعد مخفی $D$، بار محاسباتی هسته از فرم درجه دو به رابطه زیر تقلیل می‌یابد:
\begin{equation}
\label{eq:complexity_hybrid_svd_layer}
\mathcal{C}_{\text{hybrid\_svd\_layer}} = \mathcal{O}\left( B \cdot N \cdot D \cdot r + B \cdot H \cdot N \cdot d_k \right) \ll \mathcal{O}\left( B \cdot H \cdot N^2 \cdot d_k \right)
\end{equation}
این بهبود مرتبه پیچیدگی، با کاهش چشمگیر حجم تبادل تانسورها در گذرگاه‌های ارتباطی کارت گرافیک، مانع از اشباع پهنای باند حافظه (پدیده \lr{Memory-Bound}) در فرآیند استنتاج طولانی‌مدت می‌گردد.

\subsection{معماری و پیاده‌سازی}
این ماژول در فایل \lr{Codes/framework/compression\_engine/hybrid\_svd\_layer.py} پیاده‌سازی شده و تحت شناسه‌های \lr{hybrid\_svd}، \lr{hybrid\_svd\_layer} و \lr{direct\_svd\_encoder} در رجیستری فشرده‌سازی به ثبت رسیده است. اجزای اصلی معماری عبارتند از:
\begin{itemize}
    \item \textbf{کلاس \lr{HybridSVDCompressor}:} مدیریت هوشمند شناسایی لایه‌های خطی و تبدیل آن‌ها به \lr{SVDLinear} با رعایت سازگاری با معماری‌های مختلف (\lr{GPT-2, RoBERTa, DistilBERT}).
    \item \textbf{پشتیبانی از تفکیک لایه‌های توجه (\lr{Target Attention Only}):} امکان فیلتر کردن لایه‌ها برای فشرده‌سازی انحصاری پروجکشن‌های \lr{q\_proj, k\_proj, v\_proj, o\_proj, c\_attn}.
    \item \textbf{ماژول موروثی \lr{SVDSingleHead}:} حفظ پیاده‌سازی تاریخی برخط جهت سازگاری با کدهای پیشین به همراه ثبت هشدار منسوخ‌شدن (\lr{Deprecation Warning}).
    \item \textbf{محاسبه نسبت فشرده‌سازی واقعی:} ارزیابی کاهش پارامترها و محاسبه ضریب فشرده‌سازی (\lr{Compression Rate}) به صورت پویا در طول فرآیند جایگزینی.
\end{itemize}

\subsection{نتایج تجربی}
ارزیابی تجربی بنچ‌مارک فشرده‌ساز \lr{hybrid\_svd\_layer} بر روی مدل‌های مرجع در جدول \ref{tab:benchmark_hybrid_svd_layer} مستند گردیده است:

\begin{table}[H]
\centering
\small
\caption{نتایج بنچ‌مارک تجربی فشرده‌ساز \lr{hybrid\_svd\_layer} بر روی مدل‌های پایه}
\label{tab:benchmark_hybrid_svd_layer}
\begin{tabular}{lccccc}
\toprule
\textbf{مدل پایه} \lr{(Base Model)} & \textbf{دقت} \lr{(Accuracy)} & \textbf{شاخص} \lr{(F1-Score)} & \textbf{حافظه} \lr{(Peak VRAM)} & \textbf{توان} \lr{(Throughput)} & \textbf{وضعیت} \\
\midrule
\lr{DistilBERT} & \lr{77.18\%} & \lr{0.7715} & \lr{258.70 MB} & \lr{155,604.92 tok/s} & \lr{PASS} \\
\lr{RoBERTa} & \lr{50.92\%} & \lr{0.6748} & \lr{630.13 MB} & \lr{81,996.21 tok/s} & \lr{PASS\_LOW} \\
\lr{GPT-2} & \lr{51.15\%} & \lr{0.6748} & \lr{619.22 MB} & \lr{67,161.52 tok/s} & \lr{PASS\_LOW} \\
\lr{Qwen} & \lr{0.00\%} & \lr{---} & \lr{---} & \lr{---} & \lr{SKIP} \\
\bottomrule
\end{tabular}
\end{table}

نتایج بنچ‌مارک نشان می‌دهد که ماژول \lr{hybrid\_svd\_layer} با اعمال ضریب فشرده‌سازی موثر \lr{1.91x} در مدل \lr{DistilBERT}، مصرف حافظه را به شکل چشمگیری به \lr{258.70 MB} کاهش داده و توان عملیاتی بالای \lr{155,604.92 tok/s} را به همراه دقت قابل قبول \lr{77.18\%} ارائه کرده و با موفقیت وضعیت \lr{PASS} را احراز نموده است. در مدل‌های \lr{RoBERTa} و \lr{GPT-2}، با توجه به فشرده‌سازی تهاجمی بیش از ۷۰ لایه خطی با نرخ بیش از دو برابر، افت دقت در غیاب کالیبراسیون داده‌محور مشاهده شده و وضعیت \lr{PASS\_LOW} ثبت گردیده است. در مدل \lr{Qwen}، به علت عدم تطابق نام‌گذاری لایه‌ها در اسکریپت قدیمی، اجرا با خطا متوقف و به صورت \lr{SKIP} مستند شده است.

\subsection{محدودیت‌ها و تحلیل شکست}
تحلیل فنی شکست و نقاط آسیب‌پذیر این ماژول به شرح زیر است:
\begin{itemize}
    \item \textbf{فروپاشی دقت در فشرده‌سازی فراگیر بدون کالیبراسیون فعال‌سازی:} در مدل‌های \lr{RoBERTa} و \lr{GPT-2}، تجزیه مستقیم \lr{SVD} بر روی تمام ماتریس‌های خطی بدون در نظر گرفتن کانال‌های پرت (\lr{Outliers}) در فعال‌سازی‌ها، منجر به خطای بازسازی انباشته در لایه‌های متوالی و افت شدید دقت به سطح شانس تصادفی (\lr{PASS\_LOW}) شده است؛ امری که ضرورت استفاده از الگوریتم‌های آگاه از فعال‌سازی مانند \lr{ASVD} را اثبات می‌کند.
    \item \textbf{ناسازگاری مدل منسوخ \lr{SVDSingleHead}:} پیاده‌سازی قدیمی محاسبه برخط \lr{SVD} بر روی توالی فعال‌سازی‌ها علاوه بر پیچیدگی زمانی شدید در طول توالی‌های طولانی، به دلیل مقادیر منفرد نزدیک به صفر موجب ایجاد گرادیان‌های بی‌نهایت (\lr{NaN/Inf}) در گذر پسرو می‌گردد.
\end{itemize}
